\section{Types de données et champs de structure}
\label{secTypesDonnees}

\subsection{Deux familles de types}

Les documents source distinguent nettement deux familles de types, à ne pas confondre.

\begin{center}
    \begin{tabular}{|p{2.8cm}|p{6.5cm}|p{5cm}|}
        \hline
        \textbf{Préfixe de type} & \textbf{Signification} & \textbf{Exemple} \\
        \hline
        \texttt{T<Nom>} & Type maison, défini par le développeur de l'application (STRUCT ou enum). & \texttt{THmi}, \texttt{THmiInputs}, \texttt{THmiOutputs}, \texttt{TTank}, \texttt{TCtrlStatusGroup}, \texttt{TMoveAbsolute}, \texttt{TPupitre} \\
        \hline
        \texttt{T\_<Nom>} & Type de la bibliothèque cliente \texttt{uniVALplc} (Stäubli), fourni par le constructeur. & \texttt{T\_StaeubliRobot}, \texttt{T\_Status}, \texttt{T\_Command}, \texttt{T\_CommandInterface}, \texttt{T\_FromRobot}, \texttt{T\_ToRobot}, \texttt{T\_CartesianPos}, \texttt{T\_JointPos}, \texttt{T\_Trsf} \\
        \hline
    \end{tabular}
\end{center}

\subsection{Préfixes de champs de structure}

Quel que soit le type de la structure, ses champs suivent un préfixe indiquant leur type
élémentaire.

\begin{center}
    \begin{tabular}{|p{2.6cm}|p{5.4cm}|p{6.2cm}|}
        \hline
        \textbf{Préfixe} & \textbf{Type de champ} & \textbf{Exemple} \\
        \hline
        \texttt{m\_x<Nom>} & \texttt{BOOL} & \texttt{m\_xEnable}, \texttt{m\_xBusy}, \texttt{m\_xDone}, \texttt{m\_xError}, \texttt{m\_xExecute}, \texttt{m\_xActive} \\
        \hline
        \texttt{m\_udi<Nom>} & \texttt{UDINT} & \texttt{m\_udiErrorID} \\
        \hline
        \texttt{m\_t<Nom>} & \texttt{TIME} & \texttt{m\_tElapsedTime} \\
        \hline
        \texttt{m\_i<Nom>} & \texttt{INT} & \texttt{m\_iStateGmma}, \texttt{m\_iCurvSpeed\_mm\_per\_s}, \texttt{m\_iMovementID} \\
        \hline
        \texttt{m\_int<Nom>} & \texttt{INT} (variante rencontrée) & \texttt{m\_intPartNumber} \\
        \hline
        \texttt{m\_r<Nom>} & \texttt{REAL} & \texttt{m\_rVelocity}, \texttt{m\_rAcceleration}, \texttt{m\_rCartesianVel} \\
        \hline
        \texttt{m\_ui<Nom>} & \texttt{UINT} & \texttt{m\_uiCoordSystem}, \texttt{m\_uiToolNumber}, \texttt{m\_uiBufferMode} \\
        \hline
        \texttt{m\_a<N><Nom>} ou \texttt{m\_a<Nom>} & \texttt{ARRAY} (\texttt{N} = taille optionnelle) & \texttt{m\_a4qxHex7Seg}, \texttt{m\_aStankLine} (\texttt{ARRAY OF TTank}) \\
        \hline
        \texttt{m\_<nomType>} & Champ typé sur un type composé (nom du champ = nom du type en minuscule initiale) & \texttt{m\_cartesianPos} (\texttt{T\_CartesianPos}), \texttt{m\_jointPos} (\texttt{T\_JointPos}), \texttt{m\_transitionParam} (\texttt{MC\_transitionParameter}), \texttt{m\_Inputs}/\texttt{m\_Outputs} (\texttt{THmiInputs}/\texttt{THmiOutputs}) \\
        \hline
    \end{tabular}
\end{center}

\UPSTIremarque[Incohérence \texttt{m\_int} vs \texttt{m\_i}]{%
Le TP2 (annexe 3) utilise \texttt{m\_intPartNumber} pour un champ \texttt{INT}, alors que le TP3
utilise \texttt{m\_iStateGmma} pour un champ de même type. Les deux formes (\texttt{m\_int} et
\texttt{m\_i}) coexistent dans les documents source sans qu'une règle explicite ne tranche entre
elles : il s'agit probablement d'une variante orthographique du même préfixe plutôt que de deux
conventions distinctes, mais ce document ne tranche pas et rapporte les deux formes telles
qu'observées.}

\subsection{Types de transformation géométrique et de mouvement (uniVALplc)}

\begin{center}
    \begin{tabular}{|p{3.4cm}|p{5.8cm}|p{5cm}|}
        \hline
        \textbf{Type} & \textbf{Rôle} & \textbf{Champs / remarque} \\
        \hline
        \texttt{T\_CartesianPos} & Position cartésienne absolue (repère \texttt{world} ou repère utilisateur). & Utilisé par \texttt{MC\_MoveLinearAbsolute}, \texttt{MC\_MoveDirectAbsolute}, \texttt{VAL\_ShiftPoint}. \\
        \hline
        \texttt{T\_JointPos} & Position articulaire absolue (un champ \texttt{J1}..\texttt{J6} par axe, \texttt{REAL}). & Utilisé par \texttt{MC\_MoveAxisAbsolute}. Pour un robot 4 axes, \texttt{J5} et \texttt{J6} valent 0. \\
        \hline
        \texttt{T\_Trsf} & Transformation géométrique relative (offset \texttt{X, Y, Z, RX, RY, RZ}, tous \texttt{REAL}). & Utilisé avec \texttt{VAL\_ShiftPoint} pour un calcul de position relatif. \\
        \hline
        \texttt{eMC\_BUFFER\_MODE} & Enum du mode d'empilement d'une commande de mouvement (\texttt{Aborting:=0}, \texttt{Buffered:=1}, \texttt{BlendingJoint:=6}, \texttt{BlendingCartesian:=7}). & Champ \texttt{BufferMode}/\texttt{m\_uiBufferMode} des blocs de mouvement. \\
        \hline
        \texttt{eMC\_TRANSITION\_MODE} & Enum du profil de transition entre deux mouvements enchaînés (\texttt{None:=0}, \texttt{Corner distance:=3}, \texttt{Blending distances:=10}). & Champ \texttt{TransitionMode}/\texttt{m\_uiTransitionMode}. \\
        \hline
        \texttt{MC\_TransitionParameter} & Structure \texttt{leave}/\texttt{reach} (distances de sortie/reprise de trajectoire, \texttt{REAL}). & Champ \texttt{TransitionParam}/\texttt{m\_transitionParam}. \\
        \hline
        \texttt{TMoveAbsolute} / \texttt{TMoveAxisAbs} & Type maison regroupant tous les paramètres d'appel d'un bloc \texttt{MC\_Move...Absolute} (entrées, sorties, position) pour en faciliter le passage. & \texttt{TMoveAbsolute} pour les mouvements cartésiens (TP4), \texttt{TMoveAxisAbs} pour les mouvements articulaires (TP5) -- structure quasi identique, seul le nom diffère selon le TP source. \\
        \hline
    \end{tabular}
\end{center}
